\documentclass[12pt]{article}
\usepackage[utf8]{inputenc}
\usepackage[T1]{fontenc}
\usepackage{lmodern}
\usepackage{amsmath,amssymb}
\usepackage{microtype}
\usepackage{geometry}
\geometry{margin=1in}
\title{LibrIA — Cálculos, Equações e Fórmulas\\Modelos Estático, Temporal e Híbrido}
\author{Repo: LibrIA}
\date{19 de março de 2026}

\begin{document}
\sloppy
\maketitle
\section*{Notação}
- $x\in\mathbb{R}^D$: vetor de features estáticas (dimensão $D=\text{FEATURE\_DIMENSION}$).\\
- $T$: número de estimadores (árvores) quando aplicável.\\
- $K$: número de classes.\\
- $\mathbf{X}\in\mathbb{R}^{L\times F}$: sequência temporal de comprimento $L$ (sequence\_length) e $F$ features por frame (feature\_size).\\
- $p(y=k\mid\cdot)$: probabilidade estimada para a classe $k$.
\bigskip

\section*{1. Extração e Normalização de Features}
\subsection*{1.1 Modo \textit{bounding\_box} (42 features)}
Dado o conjunto de $21$ landmarks com coordenadas normalizadas $(x_i,y_i)$:
\[
x'_{i}=x_i-\min_j x_j,\qquad y'_{i}=y_i-\min_j y_j\qquad (i=1,\dots,21).
\]
O vetor de features é
\[
\mathbf{f}_{\text{bbox}} = [x'_1,y'_1,\dots,x'_{21},y'_{21}]\in\mathbb{R}^{42}.
\]

\subsection*{1.2 Modo \textit{wrist\_relative} (63 features)}
Defina o pulso (wrist) como $p_{0}=(x_0,y_0,z_0)$. Para cada landmark $i$ com $p_i=(x_i,y_i,z_i)$:
\[
r_i = p_i - p_0,\qquad d_i=\lVert r_i\rVert_2.
\]
Se $d_{\max}=\max_i d_i>0$, normaliza-se:
\[
\tilde r_i = \frac{r_i}{d_{\max}}.
\]
O vetor de features é o empilhamento:
\[
\mathbf{f}_{\text{wrist}} = [\tilde r_1^\top,\tilde r_2^\top,\dots,\tilde r_{21}^\top]\in\mathbb{R}^{63}.
\]

\subsection*{1.3 Espelhamento (mirror) usado em inferência/treino}
- Se modo = \textit{wrist\_relative}: espelha eixo $x$:
\[
\tilde r_{i,x} \leftarrow -\tilde r_{i,x}.
\]
- Se modo = \textit{bounding\_box}: espelha como $x \leftarrow 1-x$ (já normalizado no intervalo).

\section*{2. Modelo Estático (Random Forest)}
Entrada $\mathbf{x}\in\mathbb{R}^D$.

Cada árvore $t$ produz uma predição $h_t(\mathbf{x})\in\{1,\dots,K\}$. A estimativa de probabilidade por classe é:
\[
\hat p(k\mid \mathbf{x}) = \frac{1}{T}\sum_{t=1}^{T} \mathbf{1}\{h_t(\mathbf{x})=k\}.
\]
Decisão final:
\[
\hat y = \arg\max_{k}\hat p(k\mid\mathbf{x}),\qquad
\text{confiança} = \max_k \hat p(k\mid\mathbf{x}).
\]

Observação: se o objeto scikit-learn fornece $p(\cdot)$ diretamente, usa‑se essa média internamente.

\section*{3. Arquitetura CNN Embutida (Estática)}
Construção (entrada shape $(P,C)$ — p.ex. $P=21$, $C=3$):
- Camada Conv1D: filtros $f_1=16$, kernel $k=3$, ativação ReLU.
- Conv1D: $f_2=24$, kernel $3$, ReLU.
- MaxPool1D$(s=2)$.
- Conv1D: $f_3=32$, kernel $3$, ReLU.
- GlobalAveragePooling1D.
- Dense($32$), ReLU.
- Dense($K$), Softmax.

Operação de convolução 1D para um filtro $w\in\mathbb{R}^{k\times C}$ e bias $b$:
\[
z_{t} = \sum_{j=0}^{k-1} w_j\cdot x_{t+j} + b,\qquad a_t=\text{ReLU}(z_t)=\max(0,z_t).
\]
Global average pooling sobre tempo:
\[
g = \frac{1}{T'}\sum_{t=1}^{T'} a_t.
\]
Camada densa final (logits):
\[
z = W g + b,\qquad \hat p = \text{softmax}(z),\quad \text{softmax}(z)_i=\frac{e^{z_i}}{\sum_j e^{z_j}}.
\]

Função de perda usada (treino): sparse categorical crossentropy
\[
\mathcal{L} = -\frac{1}{N}\sum_{n=1}^{N}\log \big(\hat p_{y^{(n)}}(\mathbf{x}^{(n)})\big).
\]

O otimizador: Adam (parâmetros padrão do Keras, não expostos como equação aqui).

\section*{4. Arquitetura CNN Temporal (Embeddable)}
Entrada: $\mathbf{X}\in\mathbb{R}^{L\times F}$ com $L=\text{sequence\_length}$ (tipicamente 30) e $F=\text{feature\_size}$.

Camadas (conforme implementado):
- Input$(L,F)$
- Conv1D(24,3) ReLU
- Conv1D(32,3) ReLU
- MaxPool1D(2)
- Conv1D(48,3) ReLU
- GlobalAveragePooling1D
- Dense(32) ReLU
- Dense($K$) softmax

As operações seguem as mesmas fórmulas de convolução, pooling e softmax da Secção 3.

\section*{5. Arquitetura LSTM (Modelo Temporal Host)}
Entrada: $\mathbf{X}\in\mathbb{R}^{L\times F}$.

Camadas:
- LSTM(128, return\_sequences=True)
- Dropout(0.3)
- LSTM(64)
- Dropout(0.3)
- Dense(64) ReLU
- Dense($K$) softmax

Equações internas do LSTM (por passo temporal $t$; vetores em negrito):
\[
\begin{aligned}
i_t &= \sigma(W_i x_t + U_i h_{t-1} + b_i)\\
f_t &= \sigma(W_f x_t + U_f h_{t-1} + b_f)\\
o_t &= \sigma(W_o x_t + U_o h_{t-1} + b_o)\\
\tilde c_t &= \tanh(W_c x_t + U_c h_{t-1} + b_c)\\
c_t &= f_t \odot c_{t-1} + i_t \odot \tilde c_t\\
h_t &= o_t \odot \tanh(c_t).
\end{aligned}
\]
Posteriormente, $h_{L}$ (ou saída do último LSTM) é passada por Dense/ReLU e softmax para produzir $\hat p$; perda: mesma crossentropy como acima.

\section*{6. Regras de Inferência Híbrida e Arbitragem}

Existem duas camadas de regras: (A) runtime embeddable \texttt{choose\_embedded\_prediction} e (B) \texttt{PredictionMerger} (host).

(A) Escolha no runtime embedded (regra declarativa):
- Se existir previsão temporal válida e
  \[
  \text{confidence}_{\text{temporal}} \ge \tau_{temporal}
  \quad\text{e}\quad
  \text{token}_{\text{temporal}}\in\mathcal{C}_{prioridade},
  \]
  então escolher previsão temporal.
- Senão, se
  \[
  \text{confidence}_{\text{static}} \ge \tau_{static},
  \]
  escolher previsão estática.
- Senão, nenhum token é emitido.

(B) \texttt{PredictionMerger} (host) — comportamento formal:
- Parâmetros: $\mathcal{C}_{prioridade}$ (temporal priority classes), $\tau_{temporal}$, $\tau_{static}$, $t_{cooldown}$.
- Estado: $t_{\text{active\_temporal}}$ (tempo até quando predições estáticas são suprimidas).
- Ao submeter evento temporal $e$ com token $k$ e confiança $c$:
  - Se $k\notin\mathcal{C}_{prioridade}$ OR $c<\tau_{temporal}$: ignorar (não emitir).
  - Caso contrário: emitir $e$ e atualizar
    \[
    t_{\text{active\_temporal}} \leftarrow \max\big(t_{\text{active\_temporal}},\, t_{\text{end}}(e)+t_{cooldown}\big).
    \]
- Ao submeter evento estático $s$ com confiança $c$ e início $t_{\text{start}}$:
  - Se $c<\tau_{static}$: ignorar.
  - Se $t_{\text{start}}\le t_{\text{active\_temporal}}$: ignorar (bloqueado por janela temporal ativa).
  - Caso contrário: emitir $s$.
Essas regras implementam desigualdades simples e atualizações de estado com máximo.

\section*{7. Métricas e Confiança}
- Para modelos com probabilidade de saída: confiança = $\max_k \hat p(k\mid \cdot)$.\\
- Acurácia: $\frac{1}{N}\sum_{n}\mathbf{1}\{\hat y^{(n)}=y^{(n)}\}$.\\
- Durante treino registra-se \emph{training/validation loss} e \emph{accuracy} por época.

\section*{8. Observações de Implementação Relevantes}
 - Dimensões configuradas (em \texttt{config/settings.py}):
  
  \[
  \text{FEATURE\_DIMENSIONS: bounding\_box}=42,\quad wrist\_relative=63.
  \]
  Modo padrão: \texttt{wrist\_relative}. Sequence\_length típica: $L=30$.
- Quantização embedded: modelos Keras são exportados para TFLite int8 para deployment; fluxos de quantização aplicam mapeamento de ponto flutuante para representação inteira (não detalhado aqui, pois depende do calibrator e estatísticas de referência).
 - Random Forest usa parâmetros em \texttt{MODEL\_CONFIG} (ex.: $n\_estimators=100$).

\section*{9. Referências de Códigos (onde as fórmulas são aplicadas)}
 - Pré-processamento / normalização: \texttt{utils/helpers.py} (\texttt{landmarks\_to\_relative}, \texttt{landmarks\_to\_bounding\_box}).\\
 - Modelos embarcados (CNN estático/temporal): \texttt{src/model\_training/libras\_embedded\_cnn\_trainer.py}, \texttt{src/model\_training/libras\_embedded\_temporal\_cnn\_trainer.py} (arquiteturas Conv1D, pooling, GAP, Dense, Softmax).\\
 - LSTM host: \texttt{src/model\_training/libras\_lstm\_trainer.py} (camadas LSTM e configuração).\\
 - RandomForest e pipeline estático: \texttt{src/model\_training/libras\_model\_trainer.py} e artefato \texttt{model/model.pickle}.\\
 - Lógica híbrida / arbitragem: \texttt{src/inference/prediction\_merger.py} e \texttt{src/inference/libras\_embedded\_runtime.py}.

\section*{10. Conclusão}
Este documento reúne as fórmulas essenciais para estudar a arquitetura: transformação de landmarks, operações convolucionais/recorrentes, função softmax/cross‑entropy, estimativa de probabilidade do Random Forest e regras de arbitragem híbrida. Se quiser, eu:
- gero o PDF compilado deste \LaTeX{}; ou
- salvo o .tex em \texttt{docs/} do repositório e abro um PR.

\end{document}
